% Abstract content (do not add \chapter*{} here, the template handles it)
Video retrieval is a foundational operation in modern multimedia search, recommendation, and content discovery systems. Given a user query, a retrieval system must identify relevant videos from a large candidate collection and provide reliable candidates for downstream ranking, recommendation, or user interaction. As online video collections continue to grow in scale, duration, modality, and semantic complexity, video retrieval can no longer be adequately characterized as short-clip matching under a fixed benchmark protocol. Practical retrieval systems must support heterogeneous query forms, recognize evidence at different semantic granularities, construct dependable video representations from limited visual observations, and search large candidate spaces under stringent latency and throughput constraints.

This thesis studies video retrieval as a pipeline-level problem of generalization and scalability. It is organized around three connected research questions. First, how can retrieval representations be learned so that they remain effective across diverse query types, semantic conditions, and relevance patterns, rather than being specialized to a narrow set of datasets? Second, how does the density of visual observation, including temporal coverage and spatial detail, shape the corpus-side embeddings used for retrieval, and how can retrieval quality be preserved when inference must be computationally efficient? Third, how can learned retrieval signals be made searchable over large candidate collections, where relevance modeling, approximate search, graph traversal, and hardware execution are tightly coupled? By addressing these questions, the thesis connects representation learning, video embedding generation, and large-scale candidate retrieval into a coherent methodology for building video retrieval systems that are more general, more efficient, and more practical for deployment.
